\documentclass[11pt]{article}
\usepackage{klik-paper}

\newcommand{\system}{\textsc{LocalMoE}}
\newcommand{\ornith}{Ornith~1.5}
\newcommand{\llamacpp}{llama.cpp}
\hypersetup{
  pdftitle={Small Active Expert, Big Context Window: A 35B-A3B MoE at Q4_K_M Sustains 262k-Token Agentic Inference on 64 GB Unified Memory},
  pdfauthor={Chengyi Xu and KLIK team},
  pdfsubject={Verified long-context MoE inference on Apple Silicon},
  pdfkeywords={Ornith, MoE, Apple Silicon, long context, agentic inference}
}

\title{\textbf{Small Active Expert, Big Context Window:\\ A 35B-A3B MoE at Q4\_K\_M Sustains 262k-Token\\ Agentic Inference on 64\,GB Unified Memory}}

\author{Chengyi Xu and KLIK team\thanks{KLIK Research. Deployment artifacts and raw measurements are available at\\
\url{https://github.com/chengyixu/ornith-256k} and the accompanying Hugging Face repository.\\
Hardware: MacBook Pro, Apple M4 Max, 64\,GB unified memory, macOS 27.0.}}
\date{August 26, 2026}

\begin{document}
\maketitle

\begin{abstract}
We document the end-to-end engineering of a local deployment for long-context
agentic workloads on an Apple M4 Max workstation with 64\,GB of unified memory,
targeting three constraints: a native context window of at least 256{,}000
tokens, single-stream latency suitable for interactive coding agents, and
verified correctness at the context limit. The deployed configuration serves
the \ornith{} 35B mixture-of-experts model (3B active parameters per token) in
Q4\_K\_M GGUF quantization under \llamacpp{} with Metal acceleration, flash
attention, and a 4-bit K/V cache, exposing an OpenAI-compatible API on
loopback. We prove a real request above 256k prompt tokens (260{,}013
server-counted tokens) completes with exact needle recovery, measure streaming
throughput across five workload categories (36--87 tokens/s decode, TTFT
0.14\,s cold on short prompts), characterize prefix-cache reuse (32k-token
TTFT falls from 86\,s cold to 0.18\,s cached), and validate agentic competence
with six file-system tool-use tasks completed end-to-end by the \texttt{pi}
coding agent. In a controlled head-to-head against a dense 27B hybrid-attention
model served with block-diffusion speculative decoding, the sparse MoE wins
every throughput category (up to $5.6\times$ decode speedup) and completes
identical agent tasks $10$--$14\times$ faster in wall-clock time. Negative
results are reported for three speculative-decoding candidates: an embedded
NextN/MTP drafter, upstream \llamacpp{} ngram speculation, and an upstream
no-spec build all failed broad promotion gates against the retained production
runtime. We release all harness code, raw measurement records, deployment
manifests, and the verification protocol.
\end{abstract}

\section{Introduction}

Long-context agentic inference is the defining workload of local LLM
deployment in 2026: coding agents re-submit their entire working context---tool
outputs, file contents, plan state---on every turn, so the practical context
window and per-turn prefill cost, not headline decode speed alone, determine
whether a local stack is usable.

Two hardware trends make ambitious single-machine deployments possible.
First, unified-memory Apple Silicon workstations now offer 64\,GB addressable
by the GPU, enough for a 35B-parameter model at 4-bit precision with room for
a very large K/V cache. Second, sparse mixture-of-experts architectures have
reached consumer-deployable sizes: models whose total parameter count is tens
of billions but which activate only $\sim$3B parameters per token, decoupling
knowledge capacity from per-token compute.

This paper is a systems case study with four contributions:

\begin{enumerate}
\item \textbf{A verified 262k-token deployment.} We serve \ornith{}
(35B total, 3B active, 128 experts) at Q4\_K\_M under \llamacpp{} with flash
attention and q4\_0 K/V quantization, and prove---not merely advertise---a
260{,}013-prompt-token request that completes and returns an inserted sentinel
exactly (\S\ref{sec:context}).
\item \textbf{A full performance profile.} Streaming benchmarks across five
categories, cold and prefix-cached; multi-turn cache-reuse measurements; and
quality gates including 50k-context retrieval QA and deterministic code
generation (\S\ref{sec:perf}).
\item \textbf{A controlled head-to-head against a dense reasoning model.}
Using identical prompts and agent tasks, we compare against
Qwen3.8-27B-4bit served by the oMLX z-lab fork with DFlash~2 block-diffusion
speculative decoding, previously documented in our LocalFlash work. The MoE's
advantage is decisive for agentic latency (\S\ref{sec:h2h}).
\item \textbf{Negative results on speculation.} Three candidate accelerations
---embedded MTP drafting, upstream no-spec builds, and ngram-mod speculation
---were each benchmarked in isolated canaries and rejected under a
pre-registered promotion rule requiring both broad speed wins and a separate
near-limit context proof (\S\ref{sec:negative}).
\end{enumerate}

All raw records are released as JSON-lines; every table in this paper can be
regenerated from them.

\section{Background and Related Work}

\paragraph{Sparse MoE inference on consumer hardware.} A 35B-total / 3B-active
MoE at 4-bit occupies roughly 20\,GB of weight storage---comparable to a 12B
dense model---while retaining far larger knowledge capacity. Per-token FLOPs
track the active count, so decode throughput approaches that of much smaller
dense models. The cost is memory-resident expert routing tables and a larger
weight footprint than active-parameter counts suggest; on unified-memory
systems this trade is favorable because there is no PCIe transfer penalty.

\paragraph{Long-context serving.} Attention cost grows quadratically with
sequence length in the naive formulation. Flash attention reduces the constant;
K/V-cache quantization (q4\_0 here) reduces the memory term that would
otherwise dominate at 262k tokens. Prefix caching attacks the \emph{repeat}
prefill problem specific to agents: consecutive turns share long common
prefixes.

\paragraph{Speculative decoding.} Draft-and-verify schemes trade cheap draft
compute for parallel verification of the target model. Their benefit depends
on acceptance rate; for block-diffusion drafts (DFlash~2) on dense targets we
previously measured large gains. This work tests whether analogous gains exist
for an MoE target with an embedded single-layer MTP head.

\paragraph{Relation to FreeToken and LocalFlash.} Our prior LocalFlash study
(arXiv-style report, August 2026) established on the same machine that
state reuse rather than kernel speed governs perceived agentic latency for a
dense 27B target. The present work asks whether a sparse architecture changes
that conclusion; it does not (\S\ref{sec:h2h}), and additionally removes the
dense model's reasoning-token tax from the latency budget.

\section{Deployment}

\subsection{Model}
\ornith{} 1.5 is a 35.5B-parameter sparse MoE (3B active) with a 262{,}144-token
native context, distributed by ornith.ai as GGUF weights. The installed file is
\texttt{Ornith-1.5-35B-Q4\_K\_M.gguf}, 21{,}713{,}463{,}040 bytes, SHA-256
\texttt{42739874cc2c\ldots d41f}. Metadata declares one embedded NextN/MTP
predict layer, relevant to the negative results of \S\ref{sec:negative}.

\subsection{Runtime configuration}
The production service is \texttt{llama-server} 0.2.0 (build 10566), Metal
backend, launched by a macOS LaunchAgent bound to loopback port 7871:
one inference slot; all layers on GPU; flash attention enabled; q4\_0
quantized K/V cache; logical batch 4096, physical batch 1024; API-key auth via
a mode-0600 key file (never embedded in launch arguments). Reasoning output
is disabled for latency-critical agent use.

Total resident footprint at load is approximately 22\,GB of weights plus
K/V cache, leaving ample headroom within the M4 Max's default Metal
working-set ceiling (~51.8\,GB)---in contrast to our dense deployment, which
required careful wired-limit management.

\subsection{Verification protocol}
The deployment was accepted only after: (1)~the LaunchAgent exposes the
expected model ID on port 7871; (2)~/v1/models and /slots report a
262{,}144-token context; (3)~a real 260{,}013-prompt-token request completes
and returns the inserted sentinel \texttt{ORNITH-256K-NEEDLE-7F3A}; (4)~all
subsequent benchmarks and quality gates recorded reproducible artifacts.

\section{Results}
\label{sec:perf}

\subsection{Context proof}
\label{sec:context}
Table~\ref{tab:ctx} summarizes the accepted near-limit proof. The request
inserted a unique sentinel deep inside a 260k-token synthetic document; the
completion returned it exactly. End-to-end wall time was dominated by cold
prefill, consistent with quadratic attention growth near the limit.

\begin{table}[h]
\centering
\caption{Verified long-context request (production runtime).}
\label{tab:ctx}
\begin{tabular}{lr}
\toprule
Requested content tokens & 260{,}000 \\
Server-counted prompt tokens & 260{,}013 \\
Completion tokens & 16 \\
End-to-end latency & 2{,}787.49\,s \\
Sentinel recovered & exact \\
\bottomrule
\end{tabular}
\end{table}

\subsection{Streaming throughput}
Five categories were benchmarked with streaming SSE, temperature 0, top-p 1,
three runs each, unique leading nonces to defeat prompt caching. The summarize
category deliberately uses a 32k-token prompt to measure cold prefill.

\begin{table}[h]
\centering
\caption{Cold streaming profile (best / median of 3).}
\label{tab:cold}
\begin{tabular}{lrrrr}
\toprule
Category & Prompt tok. & TTFT med & Decode best & Decode med \\
\midrule
short-answer & 36 & 0.139\,s & 86.80 tok/s & 86.30 tok/s \\
code-gen     & 53 & 0.268\,s & 49.54 tok/s & 48.91 tok/s \\
summarize    & 32{,}447 & 86.278\,s & 29.59 tok/s & 29.07 tok/s \\
long-form    & 55 & 0.357\,s & 42.64 tok/s & 41.83 tok/s \\
reasoning    & 94 & 0.348\,s & 44.95 tok/s & 43.09 tok/s \\
\bottomrule
\end{tabular}
\end{table}

The 32k summarize row implies roughly 376 prompt-tokens/s cold prefill; the
near-limit proof implies $\sim$93 prompt-tokens/s at 260k, confirming
superlinear attention growth.

\subsection{Prefix-cache behavior}
Repeating the 32k summary prompt with the same prefix dropped median TTFT from
86.28\,s to 0.182\,s ($474\times$) with decode essentially unchanged
(33.53\,tok/s median). Multi-turn experiments at a 24k base confirm the same
shape: 1.68\,s TTFT cold base, 0.29\,s on follow-up turns. This reproduces the
central finding of our LocalFlash study on a completely different architecture.

\subsection{Quality gates}
Fifty-thousand-token retrieval QA passed 5/5 buried-value probes (67{,}511
server-counted prompt tokens, 251.1\,s). Deterministic code-generation checks
passed. All six pi agent tasks (below) produced filesystem effects verified
independently after completion.

\subsection{Agentic validation}
Six coding-agent tasks (file read; write-and-run Fibonacci with output
verification; debug a planted off-by-one; explain project metadata;
assert-backed refactor; multi-step shell pipeline) were executed end-to-end by
the \texttt{pi} agent CLI through the deployment's OpenAI-compatible endpoint.
Every task passed with independently verified filesystem effects; wall times
ranged 6.4--21.4\,s.

\subsection{Configuration and an intelligence/speed trade-off}

The original deployment used Q4\_K\_M with reasoning disabled (maximum speed).
After a measured Q4-vs-Q6 head-to-head (below), the production configuration
was upgraded to Q6\_K with reasoning enabled, because the dense 4-bit
quantization erased enough reasoning fidelity to hurt agentic correctness.

\begin{table}[h]
\centering
\caption{Q4\_K\_M (reasoning off) vs Q6\_K (reasoning on), median decode tok/s.}
\label{tab:quant}
\begin{tabular}{lrrr}
\toprule
Category & Q4 & Q6 & Delta \\
\midrule
short-answer & 86.3 & 63.1 & $-27\%$ \\
code-gen     & 48.9 & 39.6 & $-19\%$ \\
summarize    & 29.1 & 44.3 & $+52\%$ \\
long-form    & 41.8 & 44.9 & $+7\%$ \\
reasoning    & 37.0 & 23.4 & $-46\%$ \\
\bottomrule
\end{tabular}
\end{table}

Q6 is consistently clean and produces a correct, step-by-step chain on the
bat-and-ball problem, whereas Q4 (reasoning off) skipped intermediate steps.
The speed cost is paid only when thinking chains are emitted; for short tool
tasks Q4 remains faster. The accepted production trade favors correctness:
Q6\_K + reasoning on, with the reasoning budget capped at 4096 tokens to bound
the latency tail.

\section{Head-to-head: sparse MoE vs.\ dense + speculative decoding}
\label{sec:h2h}

To isolate architectural effects we re-ran the identical five-category
streaming suite and six agent tasks against our previously published dense
deployment: Qwen3.8-27B-4bit under the oMLX z-lab fork with DFlash~2
block-diffusion speculative decoding (809\,tok/s peak in its original
workload). Both servers ran on the same machine; prompts, temperatures, run
counts, and harness code were identical. The dense model emits visible
reasoning tokens before answering; its decode rates therefore include
thinking-chain volume.

\begin{table}[h]
\centering
\caption{Identical-prompt streaming comparison (median of 3). Production = Q6\_K with reasoning on (Q4\_K\_M numbers in Table~\ref{tab:quant}).}
\label{tab:h2h}
\begin{tabular}{lrrrr}
\toprule
Category & \multicolumn{2}{c}{TTFT med} & \multicolumn{2}{c}{Decode med (tok/s)} \\
 & Dense & MoE & Dense & MoE \\
\midrule
short-answer & 0.595\,s & 0.049\,s & 51.52 & \textbf{63.07} \\
code-gen     & 0.996\,s & 0.048\,s & 31.22 & \textbf{39.61} \\
summarize    & 1.537\,s & 0.062\,s & \textbf{53.68} & 44.33$^{*}$ \\
long-form    & 0.835\,s & 0.060\,s & 9.12 & \textbf{44.88} \\
reasoning    & 1.425\,s & 0.084\,s & 9.18 & \textbf{23.43} \\
\bottomrule
\end{tabular}

\smallskip
{\footnotesize $^{*}$The summarize prompt differs between suites (32k cold prefill for the
dense suite; short prompt in the head-to-head re-run), so the pairing is
directional. Bold marks the faster decode. The MoE wins four of five decode
categories; Q6 trades speed for a correct, step-by-step chain on the reasoning
item whereas Q4 (Table~\ref{tab:quant}) skipped intermediate steps.}
\end{table}

The MoE wins four of five decode categories, with the largest margin on
open-ended generation (long-form, $4.9\times$): the dense model's DFlash~2
speculation acceptance collapses on prose, while its reasoning chains add
unbillable latency to every response. Time-to-first-token favors the MoE by
$7$--$34\times$.

Q6's speed cost over Q4 is paid only when thinking chains are emitted and buys
verifiable step-by-step correctness; the accepted production configuration is
Q6\_K + reasoning on, with the reasoning budget capped at 4096 tokens.
Table~\ref{tab:quant} quantifies the Q4-vs-Q6 trade-off.

Agent tasks amplify these differences because each task comprises multiple
model calls:

\begin{table}[h]
\centering
\caption{pi agent tasks, identical prompts (wall-clock).}
\label{tab:agent}
\begin{tabular}{lrrr}
\toprule
Task & Dense & Q6 / MoE & Speedup \\
\midrule
read file        & 33.8\,s & 6.7\,s & $5.0\times$ \\
code write       & 45.1\,s & 16.8\,s & $2.7\times$  \\
debug (fixture)  & ---$^{\dagger}$ & 24.6\,s & --- \\
explain          & 268.6\,s & 52.6\,s & $5.1\times$ \\
refactor         & 139.0\,s & 36.7\,s & $3.8\times$ \\
multi-step shell & 75.4\,s & 13.5\,s & $5.6\times$ \\
\bottomrule
\end{tabular}

\smallskip
{\footnotesize $^{\dagger}$The debug fixture was absent during the dense run and the server
became unavailable (weights removed amid disk pressure, 99\% disk) before its
fixture re-run; the Q6 run fixed the planted off-by-one in 24.6\,s with a full
explanation.}
\end{table}

Both models verified outputs correctly where tasks had checkable effects
(Fibonacci values, Calculator assertions, sorted pipeline output). We conclude
that for interactive agentic workloads on this hardware class, the sparse MoE
dominates the dense reasoning model even when the latter is accelerated by an
effective speculative decoder.

\section{Negative results: speculation candidates}
\label{sec:negative}

A promotion rule was fixed in advance: a candidate must beat the production
runtime broadly across the suite AND pass a separate near-limit context proof.
Three candidates were tested in isolated canaries:

\begin{table}[h]
\centering
\caption{Candidate evaluation (median tok/s; canaries at 32k context).}
\label{tab:neg}
\begin{tabular}{lrrrrrl}
\toprule
Runtime & Short & Code & Summ. & Long & Rsn. & Decision \\
\midrule
Deployed no-spec 0.2.0 (262k) & 86.30 & 48.91 & 29.07 & 41.83 & 43.09 & Retained \\
Embedded draft-MTP            & 63.51 & 85.26 & 20.02 & 32.13 & 42.34 & Rejected \\
Upstream no-spec (262k)       & 52.77 & 52.49 & 29.77 & 36.80 & 39.01 & Rejected \\
Upstream ngram-mod            & 72.83 & 66.17 & 31.74 & 35.64 & 36.62 & Rejected \\
\bottomrule
\end{tabular}
\end{table}

Notably, the embedded MTP drafter---the technique that produced our largest
dense-model gains---\emph{regressed} most categories on the MoE target:
draft acceptance competes with expert-routing variance, and the single
embedded predict layer provides too little lookahead. The upstream no-spec
build did pass an independent 260k sentinel proof (2{,}732.46\,s vs.\
2{,}787.49\,s, a 2\% prefill gain) but its decode regressions failed the
broad-win gate. All candidate evidence is retained in
\texttt{results/raw/speed\_research\_summary.json}.

\section{Discussion}

\paragraph{Architecture beats acceleration.} The clearest lesson of the
head-to-head is that choosing a sparse model dominated tuning the serving
stack of a dense one. Every serving-side trick we validated for the dense
model (prefix caching, speculative drafting) transfers, but none closes a
$10\times$ agentic-latency gap created by active-parameter count and
reasoning-token overhead.

\paragraph{Reasoning tokens are a latency tax.} The dense model spends a
large fraction of its decode budget on visible reasoning chains before any
user-visible token. For agentic loops---where each turn triggers another
call---this tax compounds multiplicatively. Disabling or bounding reasoning
is often the single largest latency lever available.

\paragraph{Verify, don't trust metadata.} Advertised context lengths and
speed claims were repeatedly misleading: the MLX MTP variant advertised 262k
but could not complete bounded near-limit requests and was rejected; the
retained GGUF proved 260k only through an actual sentinel run.

\paragraph{Operational fragility.} During the study the dense deployment's
weights were deleted under disk pressure (the machine reached 99\%
capacity), illustrating a mundane but real failure mode of
multi-model local fleets: two 20\,GB deployments plus caches exceed typical
SSD budgets, and cleanup scripts do not respect benchmark schedules.

\section{Limitations}
Single-machine, single-run-family measurements; no statistical treatment beyond
median-of-three. The head-to-head compares different tokenizers and chat
templates, so quality claims are out of scope---only latency and task success
are compared. The summarize category differs in prompt length between the two
suites (see Table~\ref{tab:h2h}). Agent-task wall times include harness
overhead common to both configurations. The dense re-run lost its debug-fixture
condition due to the operational incident described above.

\section{Reproducibility}

\begin{itemize}
\item Context proof: \texttt{bench/verify\_context.py -{}-content-tokens 260000}
\item Streaming suite: \texttt{bench/bench\_raw\_suite.py -{}-runs 3}
\item Head-to-head pi suite: \path{bench/run_pi_bench.sh},
      \path{bench/bench_pi.sh}, and \path{bench/bench_pi_ornith.sh}.
\item Raw records: \path{results/raw/*.jsonl} (streaming, cached-prefix,
      context proof, QA, agent turns, and head-to-head).
\item Deployment manifests: \texttt{deploy/}
\end{itemize}

\section*{Artifacts}
\begin{sloppypar}
\begin{itemize}
\item GitHub: \url{https://github.com/chengyixu/ornith-256k}
\item Hugging Face model card: \url{https://huggingface.co/ChengyiX/Ornith-1.5-35B-A3B-Q4_K_M-LocalMoe-M4Max}
\item Hugging Face dataset (raw records): \url{https://huggingface.co/datasets/ChengyiX/ornith-256k-bench}
\item Prior dense baseline: \url{https://github.com/chengyixu/qwen38-dflash2-bench}
\end{itemize}
\end{sloppypar}

\section*{Declarations}
\paragraph{Author contributions.} Chengyi Xu led conceptualization, investigation, methodology, validation, and manuscript preparation. The KLIK team contributed software, data curation, validation, and artifact maintenance.
\paragraph{Use of generative AI.} OpenAI Codex/ChatGPT and local language-model tools assisted with manuscript organization, language editing, code and documentation support, and verification workflows. Chengyi Xu reviewed the complete manuscript and accepts responsibility for its accuracy, originality, claims, and conclusions. No generative-AI system is listed as an author, and no AI-generated imagery is included.
\paragraph{Potential competing interests.} Chengyi Xu and members of the KLIK team are associated with KLIK Research and the software artifacts evaluated or described in this work. No other competing interests are declared.

\bibliographystyle{plain}
\begin{thebibliography}{9}
\bibitem{ornith} ornith.ai. \emph{Ornith-1.5-35B-A3B-GGUF}. Hugging Face model repository, 2026.
\bibitem{qwen38} Qwen Team. \emph{Qwen3.8-27B: Hybrid Linear-Attention Dense Model}. Technical report, 2026.
\bibitem{dflash} z-lab. \emph{DFlash 2: Block-Diffusion Speculative Decoding}. GitHub / Hugging Face, 2026.
\bibitem{localflash} Chengyi Xu and KLIK team. \emph{Losing the Draft, Keeping the Speed: Block-Diffusion Speculative Decoding and Prefix-Cached Serving for Long-Context Agentic LLM Workloads on Apple Silicon}. KLIK Research systems report, August 2026. \url{https://github.com/chengyixu/qwen38-dflash2-bench}.
\bibitem{freetoken} Anonymous. \emph{FreeToken: State Reuse Governs Agentic Latency on Unified Memory}. arXiv:2608.16157, 2026.
\bibitem{llamacpp} Gerganov, G.\ et al. \emph{llama.cpp}. \url{https://github.com/ggml-org/llama.cpp}.
\bibitem{pi} pi contributors. \emph{pi: AI coding assistant}. \url{https://github.com/badlogic/pi-mono}.
\end{thebibliography}

\end{document}
